\documentclass[11pt]{article}

\usepackage[margin=1in]{geometry}
\usepackage{amsmath,amssymb,amsthm}
\usepackage{algorithm}
\usepackage{algpseudocode}
\usepackage{enumitem}
\usepackage{hyperref}
\usepackage{xcolor}
\usepackage{tikz}
\usepackage{booktabs}

\hypersetup{colorlinks=true, linkcolor=blue, urlcolor=blue, citecolor=blue}

\newtheorem{theorem}{Theorem}[section]
\newtheorem{lemma}[theorem]{Lemma}
\newtheorem{proposition}[theorem]{Proposition}
\newtheorem{corollary}[theorem]{Corollary}
\newtheorem{definition}[theorem]{Definition}
\newtheorem{example}[theorem]{Example}

\title{TOAE209/TOAH209 -- Design and Analysis of Algorithms\\[4pt]
\large Week 2: Algorithm Analysis and Asymptotic Notation}
\author{Foreign Trade University}
\date{}

\begin{document}
\maketitle
\tableofcontents

\section{Introduction}

In Week 1 we used Big-O notation informally to describe how the running time of
Gale--Shapley scales with $n$. This week we make that vocabulary precise. Following
Kleinberg \& Tardos, Chapter 2, we will:

\begin{enumerate}
    \item Define $O$, $\Omega$, and $\Theta$ rigorously, with the formal
    ``there exist constants $c, n_0$'' definitions, and prove basic properties (transitivity,
    sums, products).
    \item Build a mental hierarchy of common growth rates and learn to recognize which
    class an algorithm's running time falls into just by reading its code (counting loop
    iterations, recursive calls, etc.).
    \item Discuss \emph{polynomial time} as the standard (if imperfect) notion of
    ``efficient'', and contrast it with exponential-time brute force.
    \item Work through a full worked example -- \textbf{binary search} -- analyzing its
    running time via a recurrence relation, both directly (the iterative version) and via
    the recursion tree (the recursive version).
\end{enumerate}

This week's practical exercises (17 problems) ask you to implement and empirically verify
each of these ideas: counting operations in loops, writing ``witness checkers'' for $O$,
$\Omega$, $\Theta$, solving small recurrences by direct simulation, and analyzing classic
algorithms (binary search, heaps, dynamic arrays, Karatsuba-style divide and conquer, and a
more careful re-implementation of Gale--Shapley using arrays instead of dictionaries).

\section{Asymptotic Notation, Formally}

\subsection{Big-O: an asymptotic upper bound}

\begin{definition}[Big-O]
Let $f, g : \mathbb{N} \to \mathbb{R}^{+}$. We say $f(n) = O(g(n))$ if there exist
constants $c > 0$ and $n_0 \ge 0$ such that
\[
  f(n) \le c \cdot g(n) \qquad \text{for all } n \ge n_0.
\]
Informally: \emph{for sufficiently large $n$, $f$ grows no faster than a constant multiple
of $g$.}
\end{definition}

\begin{example}
$f(n) = 3n^2 + 100n + 7$ satisfies $f(n) = O(n^2)$. Take $c = 4$: for $n \ge n_0 = 105$,
\[
  3n^2 + 100n + 7 \le 3n^2 + n^2 + n^2 = 4n^2,
\]
because once $n \ge 105$, $100n + 7 \le n^2$ (in fact $100n \le n^2$ already holds for $n
\ge 100$, and the $+7$ is absorbed by the slack). The exact choice of $n_0$ does not matter
-- only that \emph{some} valid pair $(c, n_0)$ exists.
\end{example}

A common point of confusion: $O(g(n))$ is an \emph{upper bound} on $f$, but it does not
have to be \emph{tight}. Every function that is $O(n)$ is automatically also $O(n^2)$,
$O(n^3)$, etc. (Problem 3 in this week's exercises asks you to implement witness checkers
that make this concrete: given proposed $(c, n_0)$, verify $f(n) \le c\,g(n)$ on a sample of
$n$ values.)

\subsection{Big-Omega: an asymptotic lower bound}

\begin{definition}[Big-Omega]
$f(n) = \Omega(g(n))$ if there exist constants $c > 0$ and $n_0 \ge 0$ such that
\[
  f(n) \ge c \cdot g(n) \qquad \text{for all } n \ge n_0.
\]
Informally: \emph{for sufficiently large $n$, $f$ grows at least as fast as a constant
multiple of $g$.}
\end{definition}

Big-Omega is the mirror image of Big-O: $f(n) = \Omega(g(n))$ if and only if
$g(n) = O(f(n))$.

\subsection{Big-Theta: tight bounds}

\begin{definition}[Big-Theta]
$f(n) = \Theta(g(n))$ if $f(n) = O(g(n))$ \emph{and} $f(n) = \Omega(g(n))$. Equivalently,
there exist constants $c_1, c_2 > 0$ and $n_0$ such that
\[
  c_1 \cdot g(n) \le f(n) \le c_2 \cdot g(n) \qquad \text{for all } n \ge n_0.
\]
\end{definition}

When we say ``algorithm $A$ runs in time $\Theta(n \log n)$'', we mean the running time is
\emph{sandwiched} between two constant multiples of $n \log n$ -- this is the strongest and
most informative of the three notations, and the one we aim for whenever possible.

\subsection{Basic properties}

\begin{proposition}[Transitivity]
If $f(n) = O(g(n))$ and $g(n) = O(h(n))$, then $f(n) = O(h(n))$. The analogous statements
hold for $\Omega$ and $\Theta$.
\end{proposition}
\begin{proof}
By assumption there exist $c_1, n_1$ with $f(n) \le c_1 g(n)$ for $n \ge n_1$, and $c_2,
n_2$ with $g(n) \le c_2 h(n)$ for $n \ge n_2$. Let $n_0 = \max(n_1, n_2)$ and $c = c_1 c_2$.
For $n \ge n_0$, $f(n) \le c_1 g(n) \le c_1 c_2 h(n) = c\,h(n)$.
\end{proof}

\begin{proposition}[Sums]
If $f_1(n) = O(g(n))$ and $f_2(n) = O(g(n))$, then $(f_1+f_2)(n) = O(g(n))$. More generally,
if $f(n) = O(g(n))$ and $h(n) = O(g(n))$, then $\max(f(n), h(n)) = O(g(n))$ -- so the sum of
a constant number of terms is dominated by the asymptotically largest term, with the
others absorbed into the constant.
\end{proposition}

\begin{proposition}[Polynomials]
For a polynomial $p(n) = a_d n^d + a_{d-1} n^{d-1} + \cdots + a_0$ with $a_d > 0$,
$p(n) = \Theta(n^d)$.
\end{proposition}
\begin{proof}[Proof sketch]
For the upper bound, $p(n) \le (|a_d| + |a_{d-1}| + \cdots + |a_0|) \cdot n^d$ for $n \ge 1$
(since $n^i \le n^d$ for $i \le d$ and $n \ge 1$). For the lower bound, for $n$ large enough
the leading term $a_d n^d$ dominates the sum of the absolute values of the remaining terms,
so $p(n) \ge \frac{a_d}{2} n^d$.
\end{proof}

\section{The Hierarchy of Growth Rates}

A useful ordering of common growth rates, from slowest- to fastest-growing (each is
asymptotically dominated by the next, i.e.\ $o(\cdot)$, a \emph{strict} version of $O$):
\[
  O(1) \;\subset\; O(\log n) \;\subset\; O(\sqrt{n}) \;\subset\; O(n) \;\subset\;
  O(n \log n) \;\subset\; O(n^2) \;\subset\; O(n^3) \;\subset\; O(2^n) \;\subset\; O(n!).
\]

\begin{table}[h]
\centering
\begin{tabular}{@{}llp{6.5cm}@{}}
\toprule
Growth rate & Name & Typical example \\
\midrule
$\Theta(1)$        & constant     & array index access \\
$\Theta(\log n)$   & logarithmic  & binary search \\
$\Theta(\sqrt{n})$ & root         & some number-theoretic primality checks \\
$\Theta(n)$        & linear       & scanning a list, finding max \\
$\Theta(n \log n)$ & linearithmic & merge sort, heapsort, sorting with a heap \\
$\Theta(n^2)$      & quadratic    & insertion sort, nested-loop pairwise comparisons \\
$\Theta(n^3)$      & cubic        & naive matrix multiplication, naive triple-nested loops \\
$\Theta(2^n)$      & exponential  & brute-force subset enumeration, naive recursive Fibonacci \\
$\Theta(n!)$       & factorial    & brute-force enumeration of all permutations (e.g.\ Week 1
                                     Problem 5) \\
\bottomrule
\end{tabular}
\caption{Common running-time classes, slowest to fastest growing.}
\end{table}

\subsection{Reading growth rates off code}

The most common way a growth rate arises is by counting loop iterations or recursive
calls.

\begin{itemize}
    \item A single loop \texttt{for i in range(n)} with $O(1)$ work per iteration:
    $\Theta(n)$.
    \item A nested loop \texttt{for i in range(n): for j in range(n)}: $\Theta(n^2)$
    (Problem 2).
    \item A \emph{triangular} nested loop \texttt{for i in range(n): for j in range(i)}:
    the inner loop runs $0 + 1 + \cdots + (n-1) = \binom{n}{2} = \Theta(n^2)$ times total --
    same asymptotic class as the full nested loop, just with a smaller constant
    (Problem 2).
    \item A loop that halves its range each iteration, \texttt{while n > 1: n = n // 2}:
    $\Theta(\log n)$ iterations (Problem 5).
\end{itemize}

\subsection{Recurrences}

Many algorithms (especially recursive ones) have running times described by a
\textbf{recurrence relation}: an equation for $T(n)$ in terms of $T$ at smaller arguments.
Three recurrences recur throughout this course:

\begin{itemize}
    \item $T(n) = T(n-1) + O(1)$, $T(0) = O(1)$ $\Rightarrow$ $T(n) = \Theta(n)$ (linear
    recursion -- e.g.\ a simple countdown, Problem 4).
    \item $T(n) = T(n/2) + O(1)$, $T(1) = O(1)$ $\Rightarrow$ $T(n) = \Theta(\log n)$
    (halving recursion -- e.g.\ binary search, Problem 5).
    \item $T(n) = 2T(n/2) + O(n)$, $T(1) = O(1)$ $\Rightarrow$ $T(n) = \Theta(n \log n)$
    (divide-and-conquer with linear merge -- e.g.\ mergesort, Problem 6).
\end{itemize}

\begin{example}[Solving $T(n) = T(n/2) + 1$]
Assume $n = 2^k$. Unrolling: $T(2^k) = T(2^{k-1}) + 1 = T(2^{k-2}) + 2 = \cdots = T(2^0) + k
= T(1) + k$. Since $k = \log_2 n$, we get $T(n) = T(1) + \log_2 n = \Theta(\log n)$.
\end{example}

\begin{example}[Solving $T(n) = 2T(n/2) + n$]
Again assume $n = 2^k$. At recursion depth $i$, there are $2^i$ sub-problems, each of size
$n/2^i$, and each level contributes $2^i \cdot (n/2^i) = n$ total work (for the ``$+n$''
combine step, split proportionally across the $2^i$ sub-calls). There are $k = \log_2 n$
levels (depth $0$ through $k-1$, plus base cases), so the total work is
$n \cdot \log_2 n + \Theta(n) = \Theta(n \log n)$.
\end{example}

Problems 4--6 ask you to implement each of these three recursions \emph{exactly as
written} (so the number of recursive calls matches the formula) and verify the closed-form
solutions numerically for a range of $n$.

\section{Polynomial Time and Tractability}

\begin{definition}[Polynomial-time algorithm]
An algorithm runs in \textbf{polynomial time} if its worst-case running time is $O(n^k)$
for some constant $k$ that does not depend on the input.
\end{definition}

Kleinberg \& Tardos propose polynomial time as a useful (if imperfect) proxy for
``efficient'', for two main reasons:

\begin{enumerate}
    \item \textbf{Closure under composition.} If you call a polynomial-time subroutine a
    polynomial number of times, the overall algorithm is still polynomial-time (a
    polynomial composed with a polynomial is a polynomial). This makes polynomial-time
    algorithms \emph{compositional} -- you can build complex polynomial-time algorithms out
    of simpler polynomial-time pieces.
    \item \textbf{Sharp practical divide.} In practice, there is often a genuine
    qualitative jump between polynomial algorithms (which scale to large $n$) and
    exponential algorithms (which become infeasible for even moderately-sized $n$). An
    $O(n^3)$ algorithm on $n = 1000$ does $10^9$ operations -- seconds on a modern machine.
    A $O(2^n)$ algorithm on $n = 40$ already does over $10^{12}$ operations.
\end{enumerate}

\begin{table}[h]
\centering
\begin{tabular}{@{}lrrr@{}}
\toprule
Running time & $n=10$ & $n=30$ & $n=50$ \\
\midrule
$n$       & $10$        & $30$         & $50$ \\
$n \log n$ & $\approx 33$ & $\approx 147$ & $\approx 282$ \\
$n^2$     & $100$       & $900$        & $2{,}500$ \\
$n^3$     & $1{,}000$   & $27{,}000$   & $125{,}000$ \\
$2^n$     & $1{,}024$   & $1{,}073{,}741{,}824$ & $\approx 1.1 \times 10^{15}$ \\
$n!$      & $3{,}628{,}800$ & $\approx 2.65 \times 10^{32}$ & $\approx 3.0 \times 10^{64}$ \\
\bottomrule
\end{tabular}
\caption{Growth of running time with $n$, for fixed operation counts -- the qualitative
divide between polynomial and super-polynomial growth becomes stark by $n=30$.}
\end{table}

This is precisely why brute-force enumeration of all $n!$ matchings (Week 1, Problem 5) is
fine for $n \le 6$ but utterly infeasible for $n = 50$, while Gale--Shapley's $O(n^2)$
bound scales gracefully.

\subsection{Caveats}

Polynomial time is a useful \emph{coarse} filter, but it is not the whole story:

\begin{itemize}
    \item An $O(n^{100})$ algorithm is polynomial but useless in practice.
    \item An $O(2^{n/1000})$ algorithm is exponential but might be fine for the $n$'s that
    actually arise.
    \item Constant factors matter in practice (Theory Question 4 / Week 1 revisited):
    asymptotics are about \emph{scaling}, not absolute speed at any fixed $n$.
\end{itemize}

Nonetheless, ``is there a polynomial-time algorithm for this problem?'' remains one of the
most important questions an algorithm designer asks, and sets up the contrast (later in
the course) with $\mathsf{NP}$-hard problems for which no polynomial-time algorithm is
known.

\section{Worked Example: Binary Search}

Binary search is this week's running example because it cleanly illustrates almost every
idea above: a simple loop invariant, a $\Theta(\log n)$ recurrence, and both an iterative
and a recursive implementation that must be shown equivalent.

\subsection{Problem statement}

Given a sorted array $A[0..n-1]$ and a target value $x$, find an index $i$ such that
$A[i] = x$, or report that no such index exists.

\subsection{Iterative algorithm and loop invariant}

\begin{algorithm}[h]
\caption{\textsc{BinarySearchIterative}$(A, x)$}
\begin{algorithmic}[1]
\State $lo \gets 0$, $hi \gets n - 1$
\While{$lo \le hi$}
    \State $mid \gets \lfloor (lo + hi)/2 \rfloor$
    \If{$A[mid] = x$}
        \State \Return $mid$
    \ElsIf{$A[mid] < x$}
        \State $lo \gets mid + 1$
    \Else
        \State $hi \gets mid - 1$
    \EndIf
\EndWhile
\State \Return $-1$
\end{algorithmic}
\end{algorithm}

\begin{proposition}[Correctness via loop invariant]
At the start of every iteration, \textbf{if $x$ is present in $A$, its index lies in
$[lo, hi]$.}
\end{proposition}
\begin{proof}
\emph{Initialization.} Before the first iteration, $[lo,hi] = [0, n-1]$, the entire array,
so the invariant holds trivially.

\emph{Maintenance.} Suppose the invariant holds at the start of an iteration and $x \ne
A[mid]$. Since $A$ is sorted: if $x > A[mid]$, then $x$ (if present) must be at an index
$> mid$, so updating $lo \gets mid+1$ preserves the invariant. Symmetrically if $x <
A[mid]$, updating $hi \gets mid - 1$ preserves the invariant.

\emph{Termination.} The loop ends either by finding $x$ (return $mid$, correct by the
check), or when $lo > hi$, meaning the invariant's interval $[lo,hi]$ is empty -- by the
invariant, $x$ cannot be present, so returning $-1$ is correct.
\end{proof}

\subsection{Running-time analysis}

Each iteration either terminates or shrinks $hi - lo + 1$ (the size of the candidate range)
to \emph{at most half} its previous value (since $mid$ splits $[lo,hi]$ roughly in half,
and one half is discarded). Starting from a range of size $n$, after $k$ iterations the
range has size at most $\lceil n / 2^k \rceil$. The loop must terminate once the range size
reaches $0$, i.e.\ once $k > \log_2 n$. Hence the iterative algorithm performs
$\Theta(\log n)$ iterations, each doing $O(1)$ work, for a total running time of
$\Theta(\log n)$.

\subsection{Recursive version and recursion tree}

\begin{algorithm}[h]
\caption{\textsc{BinarySearchRecursive}$(A, x, lo, hi)$}
\begin{algorithmic}[1]
\If{$lo > hi$}
    \State \Return $-1$
\EndIf
\State $mid \gets \lfloor (lo+hi)/2 \rfloor$
\If{$A[mid] = x$}
    \State \Return $mid$
\ElsIf{$A[mid] < x$}
    \State \Return \textsc{BinarySearchRecursive}$(A, x, mid+1, hi)$
\Else
    \State \Return \textsc{BinarySearchRecursive}$(A, x, lo, mid-1)$
\EndIf
\end{algorithmic}
\end{algorithm}

This satisfies the recurrence $T(n) = T(n/2) + O(1)$ from Section 3.2, whose solution is
$T(n) = \Theta(\log n)$ -- matching the iterative analysis. The recursion ``tree'' here is
actually a single \emph{path} (each call makes exactly one recursive call), of depth
$\lceil \log_2(n+1) \rceil$. Problem 7 asks you to implement both versions and to compute
this depth directly; Problem 8 extends binary search to find the first/last occurrence of a
repeated value (still $O(\log n)$ per query), and Problem 9 applies the same
divide-and-decide idea to searching a \emph{rotated} sorted array.

\section{Other Examples of Algorithm Analysis}

\subsection{Finding the maximum and minimum: a tighter constant}

Finding the maximum of $n$ elements takes $n-1$ comparisons (one linear scan). Finding both
the max \emph{and} the min naively takes $2(n-1)$ comparisons (two scans) -- still
$\Theta(n)$, but with a worse constant. A classic trick (Problem 10) processes elements in
\emph{pairs}: for each pair, one comparison determines the local max/min candidate, then
two more comparisons update the running max and min -- $3$ comparisons per pair instead of
$4$, giving roughly $\frac{3n}{2}$ comparisons total. Both algorithms are $\Theta(n)$, but
this is a good illustration that \emph{constants matter} even within the same $\Theta$
class, and that careful counting can reveal real (if modest) improvements.

\subsection{Selection: median via sorting vs. quickselect}

Finding the median of $n$ unsorted elements by sorting first costs $\Theta(n \log n)$.
\textbf{Quickselect} (Problem 11) finds the $k$-th smallest element directly, using a
divide-and-conquer partition step (as in quicksort) but recursing into only \emph{one}
side: $T(n) = T(\alpha n) + O(n)$ for some $\alpha < 1$ in expectation, giving
$\mathbb{E}[T(n)] = \Theta(n)$ -- asymptotically faster than sorting, despite ``looking
similar'' to quicksort.

\subsection{Amortized analysis: dynamic arrays}

A \textbf{dynamic array} (Problem 15) starts with some capacity and \emph{doubles} its
capacity (copying all existing elements) whenever it becomes full. A single \texttt{append}
that triggers a resize costs $\Theta(n)$ -- but resizes happen only when the size passes
$1, 2, 4, 8, \ldots$, so the \emph{total} copying cost over $n$ appends is
$1+2+4+\cdots+2^{k-1} < 2 \cdot 2^{k-1} < 2n$, i.e.\ $O(n)$ total, or $O(1)$
\textbf{amortized} per append. This is the standard justification for treating
\texttt{list.append} in Python (and similarly, \texttt{ArrayList}/\texttt{Vector} in other
languages) as ``$O(1)$'' despite occasional expensive resizes.

\subsection{Priority queues: heap vs. sorted list}

A binary min-heap (Problem 12, building on Week 1's heapsort) supports \texttt{push} and
\texttt{pop} in $O(\log n)$, by maintaining the heap property along a single root-to-leaf
path. A naive ``always keep the list sorted'' priority queue (Problem 14) instead pays
$O(n)$ per operation (shifting elements to maintain sorted order). Problem 14 asks you to
measure this difference directly by counting ``element move'' operations for both
structures on the same sequence of pushes/pops -- a concrete $O(n \log n)$ vs.\ $O(n^2)$
comparison for $n$ insertions followed by $n$ extractions.

Problem 13 extends the heap with a \texttt{decrease\_key} operation (an \textbf{indexed
priority queue}), which is exactly the data structure needed for Dijkstra's algorithm,
covered in a later week.

\subsection{Revisiting Gale--Shapley with arrays}

Week 1's Gale--Shapley implementation used Python dictionaries and \texttt{.index()} calls
to look up preferences and ranks -- each \texttt{.index()} call is itself $O(n)$, which can
silently turn an ``$O(n^2)$'' algorithm into $O(n^3)$ if you are not careful! Problem 16
re-implements the men-proposing algorithm using only arrays (lists) indexed by integers
$0..n-1$, a \texttt{collections.deque} for the queue of free men, and a precomputed
\textbf{rank table} \texttt{women\_rank[w][m]} so that comparing two suitors is $O(1)$.
This makes the true $O(n^2)$ bound from Week 1's Lemma (at most $n^2$ proposals, each
$O(1)$ to process) explicit and verifiable.

\subsection{Correctness via induction and invariants}

Finally, Problem 17 connects this week's themes back to \emph{correctness} (not just
running time): an iterative summation maintains the loop invariant ``$total =
0+1+\cdots+i$ after processing $i$'', verified at every step with an \texttt{assert}; and
\textbf{binary exponentiation} (repeated squaring) computes $a^b$ in $O(\log b)$
multiplications instead of $O(b)$, by the recurrence
\[
  a^b = \begin{cases}
    1 & b = 0 \\
    (a^{b/2})^2 & b \text{ even} \\
    (a^{(b-1)/2})^2 \cdot a & b \text{ odd}
  \end{cases}
\]
which is the same ``halving recursion'' $T(n) = T(n/2) + O(1)$ from Section 3.2, applied to
the \emph{exponent} rather than the array size.

\section{Summary and Looking Ahead}

This week established:
\begin{itemize}
    \item The formal definitions of $O$, $\Omega$, $\Theta$, with proofs of transitivity
    and the polynomial rule.
    \item The standard hierarchy of growth rates and how to read them off code (loops,
    nested loops, recursion).
    \item Polynomial time as a (coarse but useful) definition of tractability, and why the
    polynomial/exponential divide matters in practice.
    \item A complete worked example (binary search) analyzed both iteratively (loop
    invariant + shrinking range) and recursively (recurrence $T(n) = T(n/2) + O(1)$), plus
    several further examples: max/min, quickselect, amortized analysis of dynamic arrays,
    heaps vs.\ sorted lists, and a tighter re-analysis of Gale--Shapley.
\end{itemize}

Next week, we put this analysis toolkit to use in earnest with \textbf{Divide and Conquer}
algorithms (Kleinberg \& Tardos, Chapter 5): mergesort, Karatsuba multiplication (previewed
in Week 1), and the master-theorem-style recurrences introduced informally in Section 3.2
above.

\end{document}
